\chapter{Hasil dan Pembahasan}

\section{Pembuatan Project Laravel}

Tahap pertama dalam praktikum adalah membuat aplikasi web sederhana menggunakan framework Laravel. Project dibuat menggunakan Composer dengan menjalankan perintah \texttt{composer create-project laravel/laravel}. Setelah proses instalasi selesai, project dapat dijalankan menggunakan perintah \texttt{php artisan serve}. Aplikasi kemudian dapat diakses melalui alamat lokal yang disediakan oleh Laravel.

% Screenshot 1: Proses pembuatan project Laravel
%\begin{figure}[H]
%    \centering
%    \includegraphics[width=0.7\textwidth]{2-1.png}
%    \caption{Proses pembuatan project Laravel menggunakan Composer}
%    \label{fig:2-1}
%\end{figure}

Berdasarkan Gambar 2.1, project Laravel berhasil dibuat pada direktori yang telah ditentukan. Struktur awal project secara otomatis dibuat oleh Laravel dan terdiri atas beberapa direktori seperti \texttt{app}, \texttt{config}, \texttt{database}, \texttt{public}, \texttt{resources}, \texttt{routes}, dan \texttt{tests}.

% Screenshot 2: Menjalankan Laravel dev server
%\begin{figure}[H]
%    \centering
%    \includegraphics[width=0.7\textwidth]{2-2.png}
%    \caption{Menjalankan aplikasi Laravel menggunakan Laravel development server}
%    \label{fig:2-2}
%\end{figure}

Gambar 2.2 menunjukkan bahwa aplikasi Laravel berhasil dijalankan menggunakan Laravel development server. Server lokal digunakan untuk memastikan aplikasi dapat berjalan dengan baik sebelum dilakukan proses version control dan pengunggahan ke GitHub.

\section{Pembuatan Aplikasi Web Sederhana}

Setelah project Laravel berhasil dibuat, dilakukan perubahan pada halaman utama aplikasi. Halaman tersebut digunakan untuk menampilkan informasi sederhana mengenai project praktikum, seperti judul project, identitas mahasiswa, NIM, kelas, dan deskripsi singkat aplikasi.

Perubahan halaman utama dilakukan pada file \texttt{resources/views/welcome.blade.php}. Halaman tersebut menggunakan Blade sebagai template engine yang disediakan oleh Laravel.

% Screenshot 3: Kode halaman utama
%\begin{figure}[H]
%    \centering
%    \includegraphics[width=0.7\textwidth]{2-3.png}
%    \caption{Kode halaman utama aplikasi Laravel}
%    \label{fig:2-3}
%\end{figure}

Berdasarkan Gambar 2.3, halaman utama telah disesuaikan dengan kebutuhan aplikasi sederhana. Perubahan ini menjadi salah satu bagian dari pengembangan aplikasi yang nantinya dicatat dalam riwayat commit Git.

% Screenshot 4: Tampilan halaman utama
%\begin{figure}[H]
%    \centering
%    \includegraphics[width=0.7\textwidth]{2-4.png}
%    \caption{Tampilan halaman utama aplikasi Laravel}
%    \label{fig:2-4}
%\end{figure}

Gambar 2.4 menunjukkan hasil aplikasi setelah halaman utama berhasil dikembangkan. Aplikasi dapat diakses melalui server lokal dan menampilkan informasi project yang telah ditambahkan.

Selain halaman utama, dibuat pula halaman About untuk memberikan informasi mengenai project. Route halaman About ditambahkan pada file \texttt{routes/web.php}.

% Screenshot 5: Halaman About
%\begin{figure}[H]
%    \centering
%    \includegraphics[width=0.7\textwidth]{2-5.png}
%    \caption{Tampilan halaman About aplikasi}
%    \label{fig:2-5}
%\end{figure}

Gambar 2.5 menunjukkan halaman About yang berisi informasi mengenai project praktikum. Halaman ini dapat diakses melalui route yang telah ditambahkan pada file \texttt{routes/web.php}.

\section{Inisialisasi Git dan Repository GitHub}

Setelah aplikasi berhasil dibuat, Git digunakan sebagai sistem version control untuk mengelola perubahan source code. Repository Git lokal diinisialisasi menggunakan perintah \texttt{git init}. Seluruh file project kemudian ditambahkan ke staging area menggunakan perintah \texttt{git add .}. Setelah itu dibuat commit awal menggunakan format Conventional Commits dengan pesan:

\begin{lstlisting}
chore: initialize laravel project
\end{lstlisting}

Repository baru kemudian dibuat pada organization KEPL2026 di GitHub dengan nama \texttt{evolusi-pl-24-541582-24926}. Repository dibuat dengan status \textit{Public} sesuai dengan ketentuan tugas. Repository lokal dihubungkan dengan repository GitHub menggunakan perintah \texttt{git remote add origin}, dan source code dikirimkan menggunakan perintah \texttt{git push}.

% Screenshot 6: Repository GitHub
%\begin{figure}[H]
%    \centering
%    \includegraphics[width=0.7\textwidth]{2-6.png}
%    \caption{Repository project pada GitHub}
%    \label{fig:2-6}
%\end{figure}

Gambar 2.6 menunjukkan repository \texttt{evolusi-pl-24-541582-24926} telah berhasil dibuat pada organization KEPL2026 di GitHub dengan status publik.

\section{Penerapan Conventional Commits}

Dalam proses pengembangan digunakan Conventional Commits untuk memberikan struktur yang konsisten pada riwayat perubahan. Setiap perubahan diberikan tipe commit yang sesuai dengan jenis perubahan yang dilakukan. Beberapa commit yang digunakan dalam pengembangan project adalah:

\begin{lstlisting}
chore: initialize laravel project
feat: adding wow on description
feat: add about page
feat: changing homepage content
test: add homepage feature test
ci: add github actions workflow
fix: update CI PHP version to 8.4
\end{lstlisting}

Commit dengan tipe \texttt{chore} digunakan untuk perubahan yang berkaitan dengan konfigurasi atau pekerjaan pendukung, sedangkan \texttt{feat} digunakan untuk penambahan fitur. Commit dengan tipe \texttt{test} digunakan untuk perubahan yang berkaitan dengan pengujian, \texttt{ci} digunakan untuk perubahan pada konfigurasi Continuous Integration, dan \texttt{fix} digunakan untuk perbaikan.

% Screenshot 7: Riwayat commit
%\begin{figure}[H]
%    \centering
%    \includegraphics[width=0.7\textwidth]{2-7.png}
%    \caption{Riwayat commit menggunakan Conventional Commits}
%    \label{fig:2-7}
%\end{figure}

Berdasarkan Gambar 2.7, project memiliki lebih dari lima commit dengan format Conventional Commits. Dengan demikian, setiap perubahan project dapat dilacak melalui riwayat Git secara jelas dan terstruktur.

\section{Pembuatan Branch Development dan Feature}

Tahap berikutnya adalah membuat branch \texttt{dev} dari branch \texttt{main}. Branch \texttt{dev} digunakan sebagai branch pengembangan sebelum perubahan digabungkan ke branch utama. Branch dibuat menggunakan perintah:

\begin{lstlisting}
git checkout -b dev
\end{lstlisting}

Setelah branch \texttt{dev} dibuat, dilanjutkan dengan membuat feature branch dari branch \texttt{dev} menggunakan perintah:

\begin{lstlisting}
git checkout -b feature/homepage
\end{lstlisting}

Penggunaan feature branch memungkinkan perubahan fitur dikembangkan secara terpisah dari branch \texttt{dev} dan \texttt{main}. Perubahan pada halaman utama aplikasi dikerjakan pada branch \texttt{feature/homepage}, kemudian di-commit dan dikirimkan ke GitHub.

% Screenshot 8: Branch pada repository
%\begin{figure}[H]
%    \centering
%    \includegraphics[width=0.7\textwidth]{2-8.png}
%    \caption{Branch main, dev, dan feature/homepage pada repository}
%    \label{fig:2-8}
%\end{figure}

Gambar 2.8 menunjukkan bahwa ketiga branch telah berhasil dibuat dan tersedia pada repository GitHub. Alur pengembangan yang diterapkan adalah \texttt{feature/homepage} $\rightarrow$ \texttt{dev} $\rightarrow$ \texttt{main}.

\section{Implementasi GitHub Actions}

Continuous Integration diterapkan menggunakan GitHub Actions. Workflow dibuat pada direktori \texttt{.github/workflows/ci.yml}. Workflow tersebut memiliki dua job utama, yaitu \texttt{test} dan \texttt{lint}. Job \texttt{test} digunakan untuk menjalankan pengujian aplikasi Laravel menggunakan perintah \texttt{php artisan test}, sedangkan job \texttt{lint} digunakan untuk melakukan pemeriksaan terhadap aplikasi menggunakan perintah \texttt{php artisan --version}.

Berikut adalah konfigurasi workflow yang digunakan:

\begin{lstlisting}[language=yaml]
name: Laravel CI

on:
  push:
    branches:
      - main
      - dev
      - "feature/**"
  pull_request:
    branches:
      - main
      - dev

jobs:
  test:
    runs-on: ubuntu-latest
    steps:
      - name: Checkout repository
        uses: actions/checkout@v4
      - name: Setup PHP
        uses: shivammathur/setup-php@v2
        with:
          php-version: "8.4"
          extensions: mbstring, dom, fileinfo, sqlite3
          coverage: none
      - name: Install Composer dependencies
        run: composer install --no-interaction --prefer-dist
      - name: Prepare environment
        run: |
          cp .env.example .env
          php artisan key:generate
          touch database/database.sqlite
      - name: Run Laravel tests
        run: php artisan test

  lint:
    runs-on: ubuntu-latest
    steps:
      - name: Checkout repository
        uses: actions/checkout@v4
      - name: Setup PHP
        uses: shivammathur/setup-php@v2
        with:
          php-version: "8.4"
          extensions: mbstring, dom, fileinfo, sqlite3
          coverage: none
      - name: Install Composer dependencies
        run: composer install --no-interaction --prefer-dist
      - name: Check Laravel application
        run: php artisan --version
\end{lstlisting}

% Screenshot 9: File ci.yml
%\begin{figure}[H]
%    \centering
%    \includegraphics[width=0.7\textwidth]{2-9.png}
%    \caption{Konfigurasi workflow GitHub Actions pada file ci.yml}
%    \label{fig:2-9}
%\end{figure}

Berdasarkan Gambar 2.9, file \texttt{ci.yml} telah berisi konfigurasi workflow yang terdiri dari dua job. Workflow dijalankan secara otomatis ketika terdapat perubahan pada branch yang telah ditentukan maupun ketika terdapat Pull Request.

% Screenshot 10: GitHub Actions berhasil
%\begin{figure}[H]
%    \centering
%    \includegraphics[width=0.7\textwidth]{2-10.png}
%    \caption{Hasil eksekusi GitHub Actions dengan status test dan lint berhasil}
%    \label{fig:2-10}
%\end{figure}

Gambar 2.10 menunjukkan hasil eksekusi GitHub Actions. Job \texttt{test} dan \texttt{lint} berhasil dijalankan hingga selesai dengan status berhasil. Hal tersebut menunjukkan bahwa proses Continuous Integration telah berjalan sesuai dengan konfigurasi workflow.

\section{Pull Request Feature ke Dev}

Setelah pengembangan pada branch \texttt{feature/homepage} selesai, dibuat Pull Request pertama untuk menggabungkan perubahan ke branch \texttt{dev}. Pull Request pertama memiliki alur:

\begin{lstlisting}
feature/homepage -> dev
\end{lstlisting}

Judul Pull Request yang digunakan adalah \texttt{feat: improve homepage} dengan deskripsi yang mencakup perubahan yang dilakukan serta status pengujian. Sebelum Pull Request di-merge, GitHub Actions dijalankan terlebih dahulu untuk memastikan semua pemeriksaan berhasil.

% Screenshot 11: PR feature ke dev
%\begin{figure}[H]
%    \centering
%    \includegraphics[width=0.7\textwidth]{2-11.png}
%    \caption{Pull Request dari feature/homepage ke dev}
%    \label{fig:2-11}
%\end{figure}

Gambar 2.11 menunjukkan Pull Request pertama yang mengarah dari branch \texttt{feature/homepage} menuju branch \texttt{dev}. Status pemeriksaan GitHub Actions menunjukkan bahwa job \texttt{test} dan \texttt{lint} berhasil dijalankan.

% Screenshot 12: PR feature ke dev merged
%\begin{figure}[H]
%    \centering
%    \includegraphics[width=0.7\textwidth]{2-12.png}
%    \caption{Pull Request feature/homepage ke dev berhasil di-merge}
%    \label{fig:2-12}
%\end{figure}

Berdasarkan Gambar 2.12, setelah pemeriksaan berhasil, Pull Request kemudian di-merge ke branch \texttt{dev}. Dengan demikian, perubahan dari feature branch telah masuk ke branch pengembangan.

\section{Pull Request Dev ke Main}

Setelah perubahan dari feature branch berhasil masuk ke branch \texttt{dev}, dibuat Pull Request kedua untuk menggabungkan perubahan dari \texttt{dev} ke \texttt{main}. Alur Pull Request kedua adalah:

\begin{lstlisting}
dev -> main
\end{lstlisting}

Judul Pull Request yang digunakan adalah \texttt{feat: merge development into main}. Penggunaan dua tahap Pull Request membuat proses pengembangan lebih terstruktur karena perubahan tidak langsung digabungkan ke branch utama.

% Screenshot 13: PR dev ke main
%\begin{figure}[H]
%    \centering
%    \includegraphics[width=0.7\textwidth]{2-13.png}
%    \caption{Pull Request dari dev ke main}
%    \label{fig:2-13}
%\end{figure}

Gambar 2.13 menunjukkan Pull Request kedua dari branch \texttt{dev} menuju branch \texttt{main}. Status pemeriksaan GitHub Actions juga menunjukkan bahwa seluruh job berhasil dijalankan.

% Screenshot 14: PR dev ke main merged
%\begin{figure}[H]
%    \centering
%    \includegraphics[width=0.7\textwidth]{2-14.png}
%    \caption{Pull Request dev ke main berhasil di-merge}
%    \label{fig:2-14}
%\end{figure}

Berdasarkan Gambar 2.14, setelah seluruh pemeriksaan berhasil, Pull Request dari \texttt{dev} ke \texttt{main} berhasil di-merge. Dengan demikian, seluruh perubahan dari pengembangan telah masuk ke branch utama.

\section{Branch Protection}

Branch protection diterapkan pada branch \texttt{dev} dan \texttt{main}. Konfigurasi tersebut digunakan untuk menjaga agar perubahan pada branch penting tidak dilakukan secara langsung melalui \texttt{git push}, melainkan harus melalui Pull Request.

Pada kedua branch diterapkan aturan sebagai berikut:
\begin{itemize}
    \item \textit{Require a pull request before merging} -- perubahan harus dilakukan melalui Pull Request.
    \item \textit{Require status checks to pass before merging} -- pemeriksaan GitHub Actions (\texttt{test} dan \texttt{lint}) harus berhasil sebelum perubahan dapat digabungkan.
\end{itemize}

Dengan konfigurasi tersebut, alur pengembangan yang diterapkan menjadi:

\begin{lstlisting}
feature/* -> dev -> main
\end{lstlisting}

Perubahan tidak dapat dilakukan langsung ke branch \texttt{dev} maupun \texttt{main} tanpa melalui mekanisme Pull Request dan pemeriksaan CI.

% Screenshot 15: Branch protection dev
%\begin{figure}[H]
%    \centering
%    \includegraphics[width=0.7\textwidth]{2-15.png}
%    \caption{Konfigurasi branch protection pada branch dev}
%    \label{fig:2-15}
%\end{figure}

Gambar 2.15 menunjukkan konfigurasi branch protection pada branch \texttt{dev}. Aturan tersebut membantu menjaga alur pengembangan agar perubahan dilakukan melalui mekanisme yang telah ditentukan.

% Screenshot 16: Branch protection main
%\begin{figure}[H]
%    \centering
%    \includegraphics[width=0.7\textwidth]{2-16.png}
%    \caption{Konfigurasi branch protection pada branch main}
%    \label{fig:2-16}
%\end{figure}

Berdasarkan Gambar 2.16, branch \texttt{main} memiliki aturan perlindungan yang sama dengan branch \texttt{dev}, yaitu mengharuskan perubahan dilakukan melalui Pull Request dan memenuhi pemeriksaan yang ditentukan sebelum dapat digabungkan.

\section{Penambahan Collaborator}

Tahap selanjutnya adalah menambahkan collaborator pada repository GitHub. Collaborator ditambahkan melalui pengaturan repository pada menu \textit{Settings} $\rightarrow$ \textit{Collaborators}. Collaborator diberikan permission \textit{Read} yang memungkinkan collaborator melihat isi repository tanpa memberikan akses untuk melakukan perubahan langsung terhadap source code.

% Screenshot 17: Collaborator
%\begin{figure}[H]
%    \centering
%    \includegraphics[width=0.7\textwidth]{2-17.png}
%    \caption{Penambahan collaborator dengan permission Read}
%    \label{fig:2-17}
%\end{figure}

Gambar 2.17 menunjukkan collaborator telah ditambahkan pada repository dengan hak akses \textit{Read}.

\section{Hasil Akhir Repository}

Setelah seluruh tahapan selesai dilakukan, repository memiliki alur pengembangan yang terdiri dari tiga branch utama, yaitu \texttt{feature/homepage}, \texttt{dev}, dan \texttt{main}. Pengembangan fitur dilakukan pada feature branch, kemudian perubahan digabungkan ke \texttt{dev} melalui Pull Request. Setelah melalui proses pemeriksaan, perubahan dari \texttt{dev} kemudian digabungkan ke \texttt{main} melalui Pull Request kedua.

Alur pengembangan yang diterapkan dapat diilustrasikan sebagai berikut:

\begin{lstlisting}
feature/homepage
       |
       | Pull Request #1
       v
      dev
       |
       | Pull Request #2
       v
      main
\end{lstlisting}

Riwayat commit pada repository menunjukkan seluruh perubahan yang telah dilakukan:

\begin{lstlisting}
09275e4 Merge pull request #2 from KEPL2026/dev
33502c4 Merge pull request #1 from KEPL2026/feature/homepage
391462f fix: update CI PHP version to 8.4
1cd044b ci: add github actions workflow
45193bb test: add homepage feature test
f99ceed feat: changing homepage content
b3242b2 feat: add about page
4a252ba feat: adding wow on description
99625f9 chore: initialize laravel project
\end{lstlisting}

% Screenshot 18: Branch akhir
%\begin{figure}[H]
%    \centering
%    \includegraphics[width=0.7\textwidth]{2-18.png}
%    \caption{Struktur branch akhir pada repository GitHub}
%    \label{fig:2-18}
%\end{figure}

Gambar 2.18 menunjukkan kondisi akhir repository dengan tiga branch yang tersedia, yaitu \texttt{main}, \texttt{dev}, dan \texttt{feature/homepage}.

% Screenshot 19: Kondisi akhir repository
%\begin{figure}[H]
%    \centering
%    \includegraphics[width=0.7\textwidth]{2-19.png}
%    \caption{Kondisi akhir repository dan riwayat perubahan pada GitHub}
%    \label{fig:2-19}
%\end{figure}

Berdasarkan Gambar 2.19, repository telah memiliki seluruh komponen yang diperlukan termasuk workflow GitHub Actions, branch protection, dan riwayat commit menggunakan Conventional Commits.

Berdasarkan hasil keseluruhan, project Laravel telah berhasil dikembangkan dan dikelola menggunakan Git dan GitHub. Penerapan Conventional Commits, branching, Pull Request, GitHub Actions, dan branch protection membentuk alur pengembangan yang lebih terstruktur serta sesuai dengan prinsip Continuous Integration dan praktik DevOps yang menjadi bagian dari pembelajaran.
